\section{Notation}
\begin{itemize}
    \item I use $A \cong B$ to denote that there exists an isomorphism between $A$ and $B$ in whatever the appropriate category is.
        For instance if $X$ and $Y$ are topological spaces then $X \cong Y$ means $X$ and $Y$ are homeomorphic, or if $G$ and $H$ are groups then $G \cong H$ means $G$ and $H$ are isomorphic as groups.
    \item For a homotopy $h : X \times [0,1] \to Y$, I may also write $h_t : X \to Y$ with $h_t(x) := h(x,t)$.
    \item For any set $X$, I write $c_x$ to mean the constant function at $x \in X$.
    \item Given a path $\gamma : [0,1] \to X$, I write $\Bar{\gamma}$ to mean the reversed path $\Bar{\gamma}(t) := \gamma(1-t)$.
    \item For any two paths $\alpha, \beta : [0,1] \to X$ with $\alpha(1) = \beta(0)$, I write $\beta \ast \alpha$ to indicate their concatenation:
        \[
            (\beta \ast \alpha)(t) =
            \begin{cases}
                \alpha(2t) &\text{if } 0 \leq t \leq \f12 \\
                \beta(2t-1) &\text{if } \f12 \leq t \leq 1
            \end{cases}.
        \]
        The exact paramerization is unimportant, but if it ever matters this is what I mean.
    \item Unless specified otherwise, for any group $G$ I will write $e \in G$ to be the identity in $G$.
        If there is the potential for confusion between different groups, I will use subscripts to differentiate (e.g. $e_G$ is the identity in $G$, $e_H$ is the identity in $H$).
    \item I write $\R^{d \times n}_{\rank k}$ to mean the set $\set{A \in \R^{d \times n} : \rank(A) = k}$ and set $\R^{d \times n}_* = \R^{d \times n}_{\rank d}$.
\end{itemize}